\chapter{Алгоритмы повышения ветвистости и снижения перекрытия ключей}
\label{ch:fanout}

\section{Избыточные представления ключа}
% Функции compress/decompress исторически обязательны для класса операторов
% GiST, из-за чего ключ хранится в форме, не совпадающей с внешней.
% Результат: \commit{d3a4f89d8a3}{Allow no-op GiST support functions to be omitted}.

\section{Покрывающие атрибуты}
% INCLUDE-атрибуты в GiST: \commit{f2e403803fe}{Support for INCLUDE attributes in GiST},
% в SP-GiST: \commit{09c1c6ab4bc}{Support INCLUDE'd columns in SP-GiST}.
% Неключевые данные попадают только в листья, ветвистость внутренних узлов не
% страдает, но индексное сканирование становится покрывающим.

\section{Модификация кортежа на странице}
% Существовавший порядок «удалить и дописать в конец» вызывал перемещение
% больших областей страницы и разрушал локальность.
% Результат: \commit{b1328d78f88}{Invent PageIndexTupleOverwrite}.
% Опубликовано: BDAS 2017. Числа: до трёх раз на отсортированных вставках,
% около 15\,\% на случайных — перепроверить на актуальном релизе.

\section{Внутристраничное индексирование}
% Незакрытая позиция плана 2016 г. Идея: линейный просмотр узла — это O(f)
% вызовов consistent; при росте f выигрыш от ветвистости съедается стоимостью
% просмотра. Внутристраничная структура разрывает эту связь.
% Состояние: тред 38169 (2018), WiP. Требуется довести.

\section{Функция штрафа и качество разбиения}
% Задел: ICBDA 2017. В upstream не принято — обсудить, почему, и что из этого
% следует для практики.

\section{Экспериментальная оценка}

\section{Выводы по главе}
